\subsection{Source-conditioned Eikonal factor transport}
\label{sec:witnessedtransport}

Applying the one-dimensional construction in two dimensions yields two
Cartesian compositions. Let \(\Tcal_{xy}\) apply the horizontal factor and
then the vertical factor, and let \(\Tcal_{yx}\) apply them in the reverse
order. With \(\Tcal_C=\Tcal_{xy}\) as the declared Cartesian proposal, CONV
couples the two orders through one source-conditioned Riemannian action field,
without constructing a separate raster for every source.

Let the source rectangle be
\(\Omega=[0,W-1]\times[0,H-1]\), with lattice sites \(x_p\). For channel
\(c\), define
\begin{equation}
 \begin{aligned}
 g_p^{(c)}&=(D_xY_p^{(c)},D_yY_p^{(c)})^T,\\
 G_p&=\sum_{c=1}^Cg_p^{(c)}g_p^{(c)T},
 &\omega_p&=\operatorname{tr}G_p .
 \end{aligned}
 \label{eq:witnessjet}
\end{equation}
Using the same two-jet stencils and boundary closures already present in the
Cartesian factors for \(D_x,D_y\), define the dimensionless nodal quadratic
form
\begin{equation}
 Q_p=
 \begin{cases}
 I+G_p/\omega_p,&\omega_p>0,\\
 I,&\omega_p=0,
 \end{cases}
 \qquad
 M_p=\frac{Q_p}{\det(Q_p)^{1/2}}.
 \label{eq:nodalmetric}
\end{equation}
The result is SPD with \(\det M_p=1\); scalar gradient magnitude has been
removed, leaving only the measured directional shape. For \(x\) in a source
cell, let \(b_p(x)\) be its four bilinear barycentric weights and put
\begin{equation}
 \widetilde M(x)=\sum_{p\in\operatorname{cell}(x)}b_p(x)M_p,
 \qquad
 M(x)=\frac{\widetilde M(x)}{\det(\widetilde M(x))^{1/2}}.
 \label{eq:continuousmetric}
\end{equation}
Since positive barycentric combinations preserve positive definiteness,
\eqref{eq:continuousmetric} gives one continuous determinant-one SPD field on
\(\Omega\).

For each source site \(s_a=x_a\), \(1\leq a\leq A=HW\), the action is the
intrinsic Riemannian distance
\begin{equation}
 D_a(x)=\inf_{\substack{\gamma(0)=s_a,\,\gamma(1)=x\\
                        \gamma([0,1])\subseteq\Omega}}
 \int_0^1\sqrt{\dot\gamma(t)^TM(\gamma(t))\dot\gamma(t)}\,dt.
 \label{eq:sourceaction}
\end{equation}
\begin{proposition}[Well-posed constrained action]
There exist constants \(0<\mu_-\leq\mu_+<\infty\) such that
\begin{equation}
 \mu_-I\preceq M(x)\preceq\mu_+I,\qquad x\in\Omega.
 \label{eq:actionmetricbounds}
\end{equation}
Consequently, for the length metric \(d_M\) induced by
\eqref{eq:sourceaction},
\begin{equation}
 \sqrt{\mu_-}\,\norm{x-z}_2
 \leq d_M(x,z)
 \leq \sqrt{\mu_+}\,\norm{x-z}_2 .
 \label{eq:actiondistancebounds}
\end{equation}
The constrained infimum is attained for every \(x,z\in\Omega\). Hence
\(D_a(x)=d_M(s_a,x)\) is finite and continuous, vanishes only at \(s_a\),
and is single-valued even where several minimizing curves have the same
length.
\end{proposition}

\begin{proof}
Continuity and positive definiteness of \(M\) on compact \(\Omega\) give
\eqref{eq:actionmetricbounds}. The lower inequality in
\eqref{eq:actiondistancebounds} follows by bounding the metric length of
every admissible curve below by its Euclidean length; the upper inequality
follows from the straight segment in the convex rectangle. Thus \(d_M\) is
bi-Lipschitz equivalent to the Euclidean metric, so \((\Omega,d_M)\) is a
compact length space. A compact length space is geodesic, and therefore the
infimum is attained.
\end{proof}

Equivalently, the same action is obtained as the nonnegative viscosity
solution of
\begin{equation}
 \nabla D_a(x)^TM(x)^{-1}\nabla D_a(x)=1,
 \qquad D_a(s_a)=0,
 \label{eq:eikonalpde}
\end{equation}
where paths are constrained to the closed source rectangle. In practice,
anisotropic fast marching provides a direct Cartesian evaluation of
\eqref{eq:sourceaction} \cite{mirebeau2014}.
Possible nonuniqueness of a minimizing curve introduces no branch choice in
CONV, because \eqref{eq:actionpartition} depends only on its minimum length.

Using these distances, the normalized inverse-square actions are
\begin{equation}
 w_a(x)=\frac{D_a(x)^{-2}}{\sum_bD_b(x)^{-2}},
 \qquad w_a(x)\geq0,\qquad\sum_aw_a(x)=1.
 \label{eq:actionpartition}
\end{equation}
At a source site, continuous extension of the quotient assigns unit mass to
its zero-action source site and zero mass to all others. Thus
\eqref{eq:actionpartition} is Shepard inverse-distance weighting
\cite{shepard1968} in the source-derived Riemannian distance; normalization
removes the distance units and yields a direct all-source partition without
selecting an owner.

\begin{lemma}[Kronecker action reduction]
For every pair of source sites,
\begin{equation}
 w_a(x_p)=\delta_{ap}.
 \label{eq:kroneckeraction}
\end{equation}
Consequently, for every source-indexed quantity \(F_p\),
\begin{equation}
 \sum_pw_a(x_p)F_p=F_a.
 \label{eq:nodalmomentcollapse}
\end{equation}
In particular, the action-weighted support inertia about the source site is
identically zero:
\begin{equation}
 \sum_pw_a(x_p)(x_p-s_a)(x_p-s_a)^T=0.
 \label{eq:zeroinertia}
\end{equation}
\end{lemma}

\begin{proof}
Since the Riemannian distance vanishes at its own source site and is strictly
positive between distinct lattice sites, the zero-action extension of
\eqref{eq:actionpartition} gives \eqref{eq:kroneckeraction}. Substitution then
yields \eqref{eq:nodalmomentcollapse} and \eqref{eq:zeroinertia}.
\end{proof}

The lemma removes an apparent geometric degree of freedom: any frame inferred
from \eqref{eq:zeroinertia} lies necessarily in its repeated-eigenvalue
branch. With the declared Cartesian tie \(t=e_x\), \(n=e_y\), the nonzero
nested route is therefore
\begin{equation}
 \mathcal C\!\left[
   \{\mathcal C[Y_{:,i}](v)\}_{i}
 \right](u)=\Tcal_{yx}Y(u,v),
 \label{eq:collapsedroute}
\end{equation}
independently of the source index. Consequently, a source frame bank and its
nested chart construction are algebraically redundant under the action
definitions, and \eqref{eq:collapsedroute} need be evaluated only once.

What remains is to determine how strongly the reverse order is admitted.
The current tensor already supplies the required directional energy, and its
normalized vertical component gives
\begin{equation}
 \eta_a=
 \begin{cases}
 \dfrac{(G_a)_{yy}}{\operatorname{tr}G_a},&\operatorname{tr}G_a>0,\\[5pt]
 \dfrac12,&\operatorname{tr}G_a=0.
 \end{cases}
 \label{eq:witnesscompatibility}
\end{equation}
Such normalized tensor energies are standard directional measures
\cite{bigun1987}. Since \((G_a)_{yy}\geq0\) and
\((G_a)_{yy}\leq\operatorname{tr}G_a\), equation
\eqref{eq:witnesscompatibility} gives \(0\leq\eta_a\leq1\) directly, without
an eigendecomposition, threshold, or fitted gain. The zero-current value is
one half, while exchanging the Cartesian axes sends \(\eta_a\) to
\(1-\eta_a\), exactly as it exchanges the two factor orders. The scalar
reverse-order admission is consequently
\begin{equation}
 \beta(x)=\sum_aw_a(x)\eta_a,
 \qquad 0\leq\beta(x)\leq1.
 \label{eq:actionadmission}
\end{equation}
and the complete two-dimensional CONV reconstruction is
\begin{equation}
 \boxed{
 \Tcal Y(x)=
 [1-\beta(x)]\Tcal_{xy}Y(x)+\beta(x)\Tcal_{yx}Y(x) .}
 \label{eq:witnessedblend}
\end{equation}

\begin{proposition}[Source-conditioned factor convexity]
For every target point, the two coefficients in
\eqref{eq:witnessedblend} are nonnegative and sum to one. Consequently
\(\Tcal Y(x)\) belongs to the componentwise convex hull of the two Cartesian
factor-order candidates. The operator preserves constants and reproduces
affine coordinate fields whenever the constituent factors do.
\end{proposition}

\begin{proof}
From \eqref{eq:witnesscompatibility},
\(0\leq\eta_a\leq1\); using \eqref{eq:actionpartition} then gives
\(0\leq\beta\leq\sum_aw_a=1\), and the two coefficients in
\eqref{eq:witnessedblend} sum to one. Since both constituent orders return
the same constant or affine field, their convex combination returns it as
well.
\end{proof}

\begin{proposition}[Linewise blend criterion]
Let \(\gamma:[a,b]\to\Omega\) be a \(C^1\) curve and, at every regular point,
put
\begin{equation}
 A=\Tcal_{xy}Y\circ\gamma,\qquad
 B=\Tcal_{yx}Y\circ\gamma,\qquad
 \theta=\beta\circ\gamma.
\end{equation}
Then the final reconstruction obeys the exact identity
\begin{equation}
 (\Tcal Y\circ\gamma)'
 =(1-\theta)A'+\theta B'+\theta'(B-A).
 \label{eq:blendderivative}
\end{equation}
Suppose \(sA'\geq0\) and \(sB'\geq0\) for a fixed
\(s\in\{-1,1\}\). The blend has \(s(\Tcal Y\circ\gamma)'\geq0\) exactly when
\begin{equation}
 s\theta'(B-A)
 \geq-\{(1-\theta)sA'+\theta sB'\}.
 \label{eq:blendcriterion}
\end{equation}
In particular, the readily checked bound
\begin{equation}
 |\theta'|\,|B-A|
 \leq(1-\theta)sA'+\theta sB'
 \label{eq:blendsufficient}
\end{equation}
is sufficient.
\end{proposition}

\begin{proof}
Applying the product rule to \eqref{eq:witnessedblend} gives
\eqref{eq:blendderivative}. Multiplication by \(s\) and rearrangement then
yields \eqref{eq:blendcriterion}, while \eqref{eq:blendsufficient} bounds its
left-hand term from below.
\end{proof}

\begin{proposition}[Pointwise convexity is not variation diminution]
Nonnegative pointwise blending of two monotone functions with common
endpoints does not by itself preserve monotonicity.
\end{proposition}

\begin{proof}
On \([0,1]\), take
\begin{equation}
 A(t)=t,\qquad B(t)=t+t(1-t),\qquad
 \theta(t)=\frac{1+\sin(4\pi t)}{2}.
\end{equation}
Both \(A\) and \(B\) are nondecreasing and share their endpoint values.
For \(T=(1-\theta)A+\theta B\), equation
\eqref{eq:blendderivative} gives
\begin{align}
 T'(0)&=\frac32>0,&
 T'(\tfrac34)&=\frac34-\frac{3\pi}{8}<0,\nonumber\\
 T'(1)&=\frac12>0.&&
\end{align}
Continuity therefore forces at least two derivative-sign transitions. Hence
the convex-hull statement alone cannot prove a sign-pair bound.
\end{proof}

The factorwise ringing-incapability theorem applies unconditionally to each
Cartesian composition in its declared order. The final spatially varying
blend adds no amplitude outside the componentwise interval spanned by those
two admitted values; its linewise sign topology is governed by
\eqref{eq:blendcriterion}, not by convexity alone.
